\subsection{Ograniczenia}
Przetłumaczono ograniczenia z języka naturalnego na język matematyczny
\addcontentsline{toc}{subsection}{Ograniczenia}
\begin{itemize}
\item Ograniczenie dolne wartości zmiennych decyzyjnych – wartości nie mogą być mniejsze od zera:
	\begin{equation}
	\forall{\substack{
			m \in months \\ 
			p \in products \\
			mc \in machines}} workTime[m][mc][p] >=0
	\end{equation}
	\begin{equation}
	\forall{\substack{
			m \in months \\ 
			p \in products}} produce[m][p] >=0
	\end{equation}
	\begin{equation}
	\forall{\substack{
			m \in months \\ 
			p \in products}} sell[m][p] >=0
	\end{equation}
	\begin{equation}
	\forall{\substack{
			m \in months \\ 
			p \in products}} stock[m][p] >=0
	\end{equation}
	\begin{equation}
	\forall{\substack{
			i \in scenarios \\
			m \in months \\ 
			p \in products}} lowerProfit[i][m][p] >=0
	\end{equation}

\item Ograniczenie czasowe pracy maszyn - Każda maszyna może pracować maksymalnie \textit{numberOfHoursInFactory} godzin w miesiącu, zatem łączny czas pracy wszystkich maszyn danego typu nie może przekroczyć iloczynu liczby dostępnych maszyn tego typu i czasu \textit{numberOfHoursInFactory}.
	\begin{equation}
		\forall{\substack{
			m \in months\\ 
			mc \in machines}}  \sum_{p \in products}
		(workTime[m][mc][p] <= machineCount[mc]*numberOfHoursInFactory)
	\end{equation}
    \item Ograniczenie wiążące czas pracy maszyn z produkcją - czas wykorzystania określonego typu maszyny jest równy sumie iloczynów liczby wytworzonych jednostek każdego produktu i czasu potrzebnego na obróbkę jednej jednostki tego produktu na danej maszynie:
	\begin{equation}
		\forall{\substack{
			m \in months\\ 
			mc \in machines\\
			p \in products}} workTime[m][mc][p] == produce[m][p]*timeToProduce[mc][p]
	\end{equation}
    \item Ograniczenie maksymalnej sprzedaży wynikające z pojemności rynku w danym miesiącu:
	\begin{equation}
		\forall{\substack{
			m \in months\\ 
			p \in products}} sell[m][p] == maxProductsInMonth[m][p]
	\end{equation}
	
    \item Warunki definiujące zmienną binarną przy przekroczeniu 80 procent chłonności rynku:
	\begin{equation}
		\forall{\substack{
			m \in months\\ 
			p \in products}}  sell[m][p] <= 0.8*maxProductsInMonth[m][p] + 1000000 * if80prec[m][p]
	\end{equation}
	\begin{equation}
		\forall{\substack{
			m \in months\\ 
			p \in products}} sell[m][p] >= 0.8*maxProductsInMonth[m][p] * if80prec[m][p]
	\end{equation}
	
    \item Ograniczenia linearyzujące oddziaływanie zmiennych binarnych na funkcję celu:
	\begin{equation}
		\forall{\substack{
			i \in scenarios\\			
			m \in months\\ 
			p \in products}} lowerProfit[i][m][p] <= 1000000 * if80prec[m][p]
	\end{equation}
	\begin{equation}
		\forall{\substack{
			i \in scenarios\\			
			m \in months\\ 
			p \in products}} lowerProfit[i][m][p] <= 0.2 * sell[m][p]*sellProfit[i][p]
	\end{equation}
	\begin{multline}
		\forall{\substack{
			i \in scenarios\\			
			m \in months\\ 
			p \in products}} 0.2 * sell[m][p]*sellProfit[i][p] - lowerProfit[i][m][p] +\\ 1000000 * if80prec[m][p] <= 1000000;
	\end{multline}
	
    \item Ograniczenie sprzedaży do liczby sztuk wyprodukowanych oraz dostępnych w magazynie. Dla pierwszego miesiąca ograniczenie przyjmuje formę:
	\begin{equation}
		\forall{\substack{
			m \in months\\ 
			p \in products}} sell[m][p] <= produce[m][p]+storageStart[p]
	\end{equation}
	Dla każdego następnego miesiąca:
	\begin{equation}
		\forall{\substack{
			m \in months\\ 
			p \in products}} sell[m][p] <= produce[m][p] + stock[m-1][p]
	\end{equation}
	
    \item Ograniczenie określające stan magazynu na koniec miesiąca jako różnicę między sumą produktów wyprodukowanych i dostępnych na początku miesiąca a liczbą sprzedanych jednostek. Dla pierwszego miesiąca:
	\begin{equation}
		\forall{\substack{
			m \in months\\ 
			p \in products}} stock[m][p]==(produce[m][p] + storageStart[p])-sell[m][p]
	\end{equation}
	Dla każdego następnego miesiąca:
	\begin{equation}
		\forall{\substack{
			m \in months\\ 
			p \in products}} stock[m][p]==(produce[m][p] + stock[m-1][p])-sell[m][p]
	\end{equation}

\end{itemize}
